% !TeX root = ../document_maitre.tex
%ligne qui m'a permis d'indiquer à TexStudio que ce document fait partis d'un document maître. De cette manière, j'ai facilité mon système d'écriture : en compilant un document, je compilait enfaite le document maître ce qui était plus facile pour moi. 

% Il n'y a donc pas de \begin{document} car la ligne tout en haut du document m'en dispense. Le document est donc non compilable sans document maître. 

\label{SGC?}

Les \glspl{sgc} sont aujourd'hui au \coeur des pratiques de gestion des collections et de l'activité des institutions conservatrices. Pour mesurer leur impact, il faut constater à quel point ils ont redéfini le concept de gestion des collections. De fait, ce terme désigne une toute autre réalité que celle d'origine.\newline

À l'origine, ces logiciels transposent toutes les pratiques de gestion des collections en vigueur au moment de leur invention dans un logiciel. Ainsi, ils assurent la pérennité des objets patrimoniaux et de leur histoire afin de léguer les connaissances. Les premiers logiciels comme GRIPHOS puis Micromusées ont pour but principal de répertorier chaque objet possédé, de consigner sa description et son évolution matérielle dans le temps. Finalement, ces premiers \glspl{sgc} ne déterminent pas de nouvelles pratiques.

L'explosion des moyens informatiques et leur implémentation à large échelle a conduit à la multiplication des bases de données institutionnelles. Chaque institution possède sa propre base dans laquelle elle décrit simplement ses collections avec ses propres méthodes et vocabulaires. Cela crée un paradoxe : conçus comme un atout à la normalisation des descriptions, ils ont paradoxalement renforcé le cloisonnement des inventaires. Cela devient évident avec l'invention du \textit{World Wide Web} qui rend envisageable la communication de ces inventaires numériques entre institutions et au public mais qui a besoin d'homogénéité. De plus, de nouvelles exigences juridiques sont développées afin de répondre à la rigueur scientifique requise par la recherche.

Les \glspl{sgc} permettent ainsi d'améliorer des pratiques préexistantes tout en amenant de nouveaux enjeux scientifiques et communicationnels. Plus qu'un outil de transition informatique, ils deviennents les vecteurs d'un éveil des consciences sur le besoin de changement des pratiques de gestion des collections.\newline

À partir de la décennie 1980, la gestion des collections commence à être repensée de même que les principes fondamentaux des \gls{sgc}. L'article de 1987 de David Bearman, spécialiste en stratégie d'information documentaire, \footcite{ArchivesMuseumInformatics} est fondamental dans la définition de ce que doit être un \gls{sgc}. 

Il y marque une rupture conceptuelle en définissant de nouveaux principes clés, afin de dépasser le simple rôle d'inventaire descriptif pour un rôle de  gestionnaire de l'activité globale autour des collections.\footcite[p.1]{ArchivesMuseumInformatics}.

Dans sa thèse, il distingue deux systèmes répondant à deux besoins différents : les \gls{SRI} et les \glspl{sgc}. Le premier se rapproche des systèmes contemporains de David Bearman. Il est concentré autour des \guillemet{\textit{[...] facilities that support curatorial description and information retrieval for collections management purposes [...]}}\footcite[p.2]{ArchivesMuseumInformatics},à savoir : la collecte et la recherche d'informations de description d'objets. Le second est une notion plus large centrée sur une gestion du cycle de vie des objets culturels. 5 notions sont fondamentale : le temps, l'espace, les fonds financiers, le personnel et les fonds de collections\footcite[p.3]{ArchivesMuseumInformatics}. C'est le fondement de nos systèmes actuels.

David Bearman soulève un paradoxe : les \gls{sgc} de première génération n'incluent pas de dimensions managériales concrètes mais uniquement du stockage et de la recherche d'informations statiques. Il y a donc un écart entre les besoins croissants des institutions souhaitant établir des stratégies d'utilisation de ces objets et ce que ces outils permettent. En somme, un véritable \gls{sgc} ne doit pas seulement consigner des connaissances descriptive, mais piloter les activités et répondre aux enjeux stratégiques des collections.\newline

Ainsi, un \gls{sgc} doit gérer des données sur le cycle de vie des objets patrimoniaux\footcite[p.4-11]{ArchivesMuseumInformatics} et des données centrées sur la gestion de ces derniers\footcite[p.12-22]{ArchivesMuseumInformatics}.

Pour le premier point, il s'agit de consigner des données descriptives et historiques de l'objet. Mais également tout ce qui va relever de son intégration dans les collections. Ainsi, on y consigne tout ce qui concerne la vie d'un objet. À savoir ses conditions d'entrées dans les collections, les mouvements successifs, et toute activité concernant l'objet.

Pour le second point, il comprend la création, la catégorisation et la réutilisation de toutes les données produites au sujet de l'activité d'un objet patrimonial au sein des collections. C'est-à-dire qu'on veut consigner et décrire avec la même précision tout ce qui gravite autour de la gestion des collections : les événements, les agents de l'institution, les autorités, c'est-à-dire les personnes, les ressoucrs et actions liées aux objets. Chacun de ces éléments doit être décrit, catégorisé, défini dans ses fonctions et limites et reliés à des objets de la collection afin d'en nourrir le cycle de vie. 

il existe ainsi un lien de dépendance important et à double sens. Le fonctionnement d'une institution façonne sa gestion des collections, et ce constat et valable à l'inverse.\newline

Sans aller plus loin dans l'article de David Bearman, on comprend alors que l'objectif d'un \gls{sgc} est d'être un outil stratégique et de pilotage des activités et politiques institutionnelles. La gestion des collections étant le \coeur de ces institutions, elle en influence nécessairement les activités institutionnelles et sa capacité d'action.

C'est pour cela que ces systèmes sont très réglementés par des cahiers des charges précis observant des conventions internationales. C'est le cas de la norme SPECTRUM \footcite{Spectrum} du \textit{Collections Trust} définissant des procédures fondamentales sur : l'entrée, l'acquisition, la localisation, l'inventaire périodique, le catalogage, la sortie, les prêts entrants, les prêts sortants et les mouvements, la conservation, la gestion des risques, l'utilisation des collections et le déclassement des objets patrimoniaux. Tout \gls{sgc} doit pouvoir les documenter, les renseigner et les tracer dans le respect de l'\gls{intero}. 

Cette norme concerne aussi bien la gestion des objets numériques, que celle des objets physiques. C'est un point important qui fait du \gls{sgc} l'outil de gestion des collections aussi bien physiques que numériques. Cela amène donc à créer un outil multi-facettes afin de correspondre à différentes réalités : celle du numérique et celle du physique. Ainsi, cet outil doit répondre aux enjeux stratégiques institutionnels mais aussi de gestion de fonds doublement hétérogènes.\newline

Ajoutons à cela que les institutions publiques ne sont pas les seules concernées : le secteur privé, notamment le luxe, s'empare désormais de ces logiciels. Pour ces entreprises, la valorisation du patrimoine est devenue un levier stratégique au service de l'image de marque. Les collections y sont envisagées comme des outils commerciaux à part entière. La quête de productivité et de rendement justifie l'adoption de \glspl{sgc} capables d'optimiser le pilotage de leurs activités.\newline

Finalement, aujourd'hui un \gls{sgc} est un outil informatique d'une grande complexité. Devenus le \coeur de la gestion des collections dès les années 1970, ils sont devenus les points centraux des institutions conservatrices. Ils assurent la gestion des collections, des stratégies patrimoniales et institutionnelles ainsi que toutes les activités et événements liés à ces deux points. 

Pour cette raison, les \gls{sgc} doivent assurer la documentation de l'objet, de sa vie au sein des collections, de son histoire mais aussi de la vie institutionnelle. Au-delà de ces aspects, les \gls{sgc} sont au \coeur de politiques internationales d'harmonisation des pratiques de gestion des collections et des données produites au sein des activités concernées. Pour cette raison, ce sont des logiciels très contrôlés avec des cahiers des charges précis et complets qu'il convient de respecter.\footnote{en témoigne la liste dressée par l'organisme \textit{Collection trusts} des logiciels approuvés : \url{https://collectionstrust.org.uk/software/}}

Cela en fait des logiciels complexes et complets qui doivent s'adapter à différentes collections, différents volumes de données et différents buts notamment avec l'arrivée du secteur privé dans le domaine de la gestion des collections. 

Ainsi, un \gls{sgc} doit affronter un grand nombre de difficultés. C'est ce que nous allons aborder ensuite. 